\documentclass[11pt,a4paper]{article}
\usepackage[T1]{fontenc}
\usepackage[utf8]{inputenc}
\usepackage[french]{babel}
\usepackage[margin=2.4cm]{geometry}
\usepackage{xcolor}
\usepackage{mdframed}
\usepackage{enumitem}
\usepackage{csquotes}

\definecolor{avant}{RGB}{150,40,40}
\definecolor{apres}{RGB}{30,110,70}
\definecolor{hdr}{RGB}{22,50,92}

\newmdenv[linecolor=avant!40,backgroundcolor=avant!4,linewidth=0.8pt,
  roundcorner=4pt,innertopmargin=6pt,innerbottommargin=6pt,
  innerleftmargin=10pt,innerrightmargin=10pt,skipabove=6pt,skipbelow=4pt]{avantbox}
\newmdenv[linecolor=apres!50,backgroundcolor=apres!5,linewidth=0.8pt,
  roundcorner=4pt,innertopmargin=6pt,innerbottommargin=6pt,
  innerleftmargin=10pt,innerrightmargin=10pt,skipabove=4pt,skipbelow=6pt]{apresbox}

\newcommand{\av}{\textbf{\textcolor{avant}{AVANT}}\quad}
\newcommand{\ap}{\textbf{\textcolor{apres}{APRÈS}}\quad}
\newcommand{\pourquoi}[1]{\par\smallskip\noindent\textit{\textbf{Pourquoi :} #1}\par\smallskip}

\title{Propositions de réécriture --- réduction de l'empreinte IA\\\large Fichier \texttt{main.tex}}
\author{Auditeur automatique (Claude Code)}
\date{\today}

\begin{document}
\maketitle

\noindent Ce document propose, passage par passage, une version réécrite des extraits les
plus marqués \enquote{IA} dans \texttt{main.tex}. Objectif : varier le rythme, supprimer
les formules génériques, ajouter de la voix personnelle et des choix assumés --- \emph{sans
toucher au fond technique}, qui est solide. Les réécritures sont des \textbf{points de
départ} à adapter à votre voix réelle ; ne recopiez pas mot pour mot.

\bigskip
\section*{1. Introduction --- ouverture}
\addcontentsline{toc}{section}{1. Introduction}
\noindent\textit{Localisation :} \texttt{\textbackslash section\{Introduction\}}, 1\textsuperscript{er} paragraphe.

\begin{avantbox}
\av La donnée s'est imposée comme l'un des actifs stratégiques majeurs des organisations
contemporaines. Derrière chaque tableau de bord, chaque modèle prédictif et chaque service
numérique se trouve une chaîne d'ingénierie qui transforme des flux bruts, massifs et
hétérogènes en information fiable, exploitable et créatrice de valeur.
\end{avantbox}
\begin{apresbox}
\ap Pendant deux ans d'alternance chez Geismar, j'ai passé le plus clair de mon temps sur
un problème simple à énoncer et difficile à résoudre : faire parler ensemble des données
qui ne se parlent pas. ERP M3, Business Central, vieilles bases Access --- chaque filiale
a son format, son rythme, ses trous. Le métier de data engineer, tel que je l'ai pratiqué,
consiste justement à construire la chaîne qui transforme ce désordre en chiffres fiables.
\end{apresbox}
\pourquoi{\enquote{s'est imposée comme}, \enquote{actifs stratégiques}, \enquote{créatrice
de valeur} sont des marqueurs d'ouverture génériques. La version réécrite ancre l'intro
dans votre expérience réelle (déjà décrite plus loin dans le rapport) et donne une voix.}

\section*{2. Remerciements}
\noindent\textit{Localisation :} \texttt{\textbackslash section*\{Remerciements\}}.

\begin{avantbox}
\av Cet environnement a joué un rôle déterminant dans mon développement professionnel en me
permettant d'acquérir de nouvelles compétences. Je tiens à souligner combien les
opportunités d'expérimentation qui m'ont été offertes ont été d'une valeur inestimable.
\end{avantbox}
\begin{apresbox}
\ap Chez Geismar, on m'a laissé tester des choses : migrer pour de vrai des données vers
Snowflake, casser puis réparer des pipelines, défendre mes choix techniques. C'est cette
liberté d'expérimenter, plus que tout le reste, qui m'a fait progresser.
\end{apresbox}
\pourquoi{\enquote{joué un rôle déterminant}, \enquote{valeur inestimable} sont des
formules toutes faites. Un exemple concret (\enquote{migrer vers Snowflake},
\enquote{casser puis réparer}) sonne sincère et humain. \textbf{À corriger aussi :}
\enquote{au intervenants} $\rightarrow$ \enquote{aux intervenants} ;
\enquote{La richesse\,\dots\ ont} $\rightarrow$ \enquote{a} ; \enquote{essentiels de ma
progression} $\rightarrow$ \enquote{essentiels à ma progression}.}

\section*{3. Bloc 1 --- paragraphe « Valeur du projet dans le rapport »}
\noindent\textit{Localisation :} Bloc 1, \texttt{\textbackslash paragraph\{Valeur du projet dans le rapport.\}}

\begin{avantbox}
\av Ce projet est particulièrement intéressant pour ce rapport, car il constitue un cas
concret d'architecture de données de bout en bout. [\dots] En d'autres termes, il ne s'agit
pas d'un simple exercice d'import de données, mais bien d'une plateforme cohérente où chaque
couche répond à un rôle précis [\dots]
\end{avantbox}
\begin{apresbox}
\ap Si je retiens ce projet pour le Bloc 1, c'est qu'il m'a obligé à traiter toute la
chaîne, pas juste un maillon : aller chercher six sources qui n'avaient ni le même format
ni le même rythme, les fiabiliser, puis les exposer via une API utilisable. C'est là que
j'ai vraiment compris ce que \enquote{Bronze / Silver / Gold} veut dire en pratique --- et
pourquoi sauter une couche se paie plus tard.
\end{apresbox}
\pourquoi{Le paragraphe d'origine est un \enquote{méta-commentaire} qui justifie le choix
sans rien apporter de neuf (\enquote{particulièrement intéressant}, \enquote{en d'autres
termes}). La réécriture remplace la justification abstraite par un apprentissage concret.}

\section*{4. Bloc 1 --- clôtures de justification répétées}
\noindent\textit{Localisation :} sous-section \enquote{Justification des compétences du bloc 1}
(plusieurs sous-parties se terminent par la même tournure).

\begin{avantbox}
\av Cette partie répond donc à la compétence attendue, car il ne s'agissait pas seulement
de\,\dots\quad(\textit{puis, plus loin :}) Cette conception répond à la compétence visée,
car\,\dots\quad(\textit{puis :}) Cette API justifie pleinement la compétence attendue\,\dots
\end{avantbox}
\begin{apresbox}
\ap (Varier les clôtures, p.\,ex.)\\
--- \enquote{Concrètement, ce sont la modélisation en étoile et l'\textit{upsert} MySQL qui
matérialisent la compétence \enquote{base relationnelle}.}\\
--- \enquote{Le choix DynamoDB se défend par les patrons d'accès, pas par la mode NoSQL :
nous sommes partis des requêtes attendues.}\\
--- \enquote{L'API n'est pas un ajout cosmétique : filtres, pagination et auth JWT en font
un vrai service consommable.}
\end{apresbox}
\pourquoi{La répétition mécanique de \enquote{répond à la compétence attendue} est un
signal de remplissage. Conclure chaque sous-partie par la \emph{preuve technique précise}
est à la fois plus humain et plus convaincant pour un jury.}

\section*{5. Bloc 2 --- phrase à rallonge et coquilles}
\noindent\textit{Localisation :} Bloc 2, \enquote{Architecture distribuée avec Spark} et
\enquote{Couche Bronze}.

\begin{avantbox}
\av L'ensemble du projet repose sur une architecture autour de Spark, nous avons installé
spark en local sur une machine puissante, le point d'entrée le cluster se fait via
\texttt{spark://localhost:7077}. [\dots] partitionnement des vols par Year, Month et Day
\textbf{cat} il améliore les performances\,\dots
\end{avantbox}
\begin{apresbox}
\ap Le projet tourne sur un cluster Spark installé en local sur une machine dédiée ; le
point d'entrée est \texttt{spark://localhost:7077}. Tous les scripts (ingestion, streaming,
nettoyage, agrégation, ML) partagent la même \texttt{SparkSession}, où nous fixons la
mémoire et le nombre de cœurs. Nous partitionnons les vols par année, mois et jour :
\textbf{car} cela limite les lectures à la période utile au lieu de relire tout le jeu de
données.
\end{apresbox}
\pourquoi{Phrase d'origine en \enquote{run-on} (plusieurs idées sans ponctuation forte) et
deux coquilles : \enquote{le point d'entrée le cluster} et \enquote{cat} $\rightarrow$
\enquote{car}. La réécriture découpe le rythme \emph{tout en gardant} votre détail concret
(\texttt{:7077}, mémoire/cœurs) qui est un excellent marqueur humain.}

\section*{6. Bloc 4 --- Architecture / cycle de vie}
\noindent\textit{Localisation :} Bloc 4, fin de \enquote{Architecture et cycle de vie du modèle}.

\begin{avantbox}
\av Ce découpage modulaire est un choix d'ingénierie [\dots] Il est impératif de séparer
les responsabilités pour que le projet ML puisse être maintenu.
\end{avantbox}
\begin{apresbox}
\ap Concrètement, ce découpage nous a déjà servi : nous avons réentraîné les modèles deux
fois sans toucher à l'API, et fait évoluer le dashboard sans relancer un seul
entraînement. C'est tout l'intérêt de garder préparation, entraînement, évaluation et
serving étanches.
\end{apresbox}
\pourquoi{\enquote{Il est impératif de\,\dots} est une formule sentencieuse générique
(burstiness de la section : 0,12, très uniforme). Donner un \emph{effet vécu}
(\enquote{réentraîné deux fois sans toucher à l'API}) casse le rythme et prouve le bénéfice.}

\section*{7. Synthèses de bloc (Blocs 3 et 4)}
\noindent\textit{Localisation :} débuts de \enquote{Synthèse --- compétences démontrées}.

\begin{avantbox}
\av Le \texttt{projet cloud} couvre l'intégralité du périmètre du Bloc 3 [\dots]
(\textit{et} :) Le \texttt{projet data science} couvre l'intégralité du périmètre du Bloc 4\,\dots
\end{avantbox}
\begin{apresbox}
\ap (Bloc 3) Au final, ce projet m'a fait toucher les quatre faces du Bloc 3 sur un même
cas : stocker (S3/DynamoDB), traiter (ingestion + agrégats), sécuriser (IAM, chiffrement)
et optimiser (GSI O(1), serverless). Le tableau ci-dessous relie chaque compétence à sa
preuve.
\end{apresbox}
\pourquoi{\enquote{couvre l'intégralité du périmètre} ouvre les deux synthèses à
l'identique : tournure de clôture stéréotypée. Rappeler les 4 axes concrets en une phrase
est plus informatif et personnel.}

\section*{8. Zones à rédiger (pas une réécriture : un squelette)}
Ces parties sont \textbf{vides} (texte de consigne). Voici une trame à remplir avec vos
éléments réels --- ne pas laisser de texte générique.

\subsection*{Bloc 5 --- Gouvernance des données}
\begin{itemize}[leftmargin=1.3em]
  \item \textbf{Contexte réel Geismar} : multiplicité ERP (M3, Business Central, Access),
  besoin d'harmonisation groupe, migration Snowflake --- un terrain de gouvernance concret.
  \item \textbf{Gouvernance} : qui possède la donnée financière / stocks ? rôles, data owners,
  catalogue, conventions de nommage Bronze/Silver/Gold.
  \item \textbf{Conformité} : RGPD sur les données RH / clients, rétention, anonymisation.
  \item \textbf{Qualité} : votre \texttt{data\_quality\_score} (Bloc 1) est une vraie preuve
  d'audit qualité --- réutilisez-le ici.
  \item \textbf{Sécurité \& accès} : IAM (Bloc 3), rôles Snowflake, moindre privilège.
  \item \textbf{Stratégie / communication} : Power BI comme vecteur de diffusion de la culture data.
\end{itemize}

\subsection*{Conclusion}
\begin{itemize}[leftmargin=1.3em]
  \item Fil conducteur : des cinq projets, du stockage (B1) à la gouvernance (B5).
  \item Progression sur deux ans : ce que vous saviez faire avant / après.
  \item Recul critique : un choix que vous referiez autrement (p.\,ex. validation k-fold,
  SMOTE non implémenté, ingestion EventBridge restée théorique).
  \item Ouverture vers la soutenance et l'e-portfolio.
\end{itemize}

\subsection*{Corrections de gabarit (rapides)}
\begin{itemize}[leftmargin=1.3em]
  \item \texttt{\textbackslash author\{Tonio\}} $\rightarrow$ vos deux noms.
  \item Page de titre : \enquote{Encadrant : A completer}, \enquote{Entreprise : A completer}.
  \item Supprimer le paragraphe \enquote{Analyse attendue} du Bloc 2.
  \item Accents des titres : \enquote{developper}, \enquote{strategie}, \enquote{methodes},
  \enquote{creatrices}, \enquote{deployer}.
  \item \enquote{projet datas cience project} (intro) $\rightarrow$ \enquote{data science project}.
  \item \enquote{decapteurs} (Bloc 4 EDA) $\rightarrow$ \enquote{des capteurs}.
\end{itemize}

\end{document}
